\appendix

\section{Full tables}\label{app:full}

Every number in this appendix is generated by the analysis script from the three result
artifacts the experiments write; none is typed by hand. Point estimates here are the same
values reported in the main text, so a reader can check one against the other.

\subsection{Recovery by failure class and policy}

\begin{table}[!ht]
\centering\scriptsize
\setlength{\tabcolsep}{4pt}
\caption{Recovery rate (\% of harmful pairs) by failure class with 95\% cluster-bootstrap intervals over task templates. \emph{Det.} is the rule monitor's event recall for that class. The unrecoverable class is reported as attempted repair; containment is scored separately.}
\label{tab:appendix_class}
\begin{adjustbox}{max width=\textwidth}
\begin{tabular}{@{}lrccccr@{}}
\toprule
Failure class & $n$ & No recovery & Retry only & Rule-based & Model-based & Det. \\
\midrule
invalid\_arguments & 169 & 0.0 [0.0, 0.0] & 0.0 [0.0, 0.0] & 95.9 [94.6, 98.2] & 73.4 [73.2, 73.7] & 100.0 \\
stale\_state & 76 & 0.0 [0.0, 0.0] & 0.0 [0.0, 0.0] & 76.3 [65.0, 89.3] & 68.4 [60.0, 75.0] & 86.8 \\
contradictory\_evidence & 158 & 0.0 [0.0, 0.0] & 0.0 [0.0, 0.0] & 96.2 [92.7, 100.0] & 96.2 [92.7, 100.0] & 100.0 \\
planning\_error & 110 & 0.0 [0.0, 0.0] & 0.0 [0.0, 0.0] & 40.9 [38.5, 42.9] & 41.8 [41.0, 42.9] & 80.0 \\
communication\_failure & 45 & 0.0 [0.0, 0.0] & 0.0 [0.0, 0.0] & 82.2 [73.3, 93.3] & 82.2 [73.3, 93.3] & 100.0 \\
objective\_drift & 134 & 0.0 [0.0, 0.0] & 0.0 [0.0, 0.0] & 52.2 [44.2, 57.8] & 23.1 [8.9, 30.4] & 97.8 \\
unrecoverable\_external\_error & 169 & 0.0 [0.0, 0.0] & 0.0 [0.0, 0.0] & 0.0 [0.0, 0.0] & 0.0 [0.0, 0.0] & 100.0 \\
\bottomrule
\end{tabular}
\end{adjustbox}
\end{table}


\subsection{Recovery by environment family}

\begin{table}[!ht]
\centering\scriptsize
\setlength{\tabcolsep}{4pt}
\caption{Recovery rate (\%) by environment family and policy with 95\% cluster-bootstrap intervals. The accounting family is hardest for every policy; the gap between rule-based and no recovery is largest where state is auditable.}
\label{tab:appendix_family}
\begin{tabular}{@{}lrcccc@{}}
\toprule
Family & $n$ & No recovery & Retry only & Rule-based & Model-based \\
\midrule
inventory & 217 & 0.0 [0.0, 0.0] & 0.0 [0.0, 0.0] & 69.1 [65.8, 73.6] & 64.1 [61.6, 68.1] \\
migration & 196 & 0.0 [0.0, 0.0] & 0.0 [0.0, 0.0] & 77.0 [75.0, 78.6] & 59.7 [56.7, 62.9] \\
procurement & 285 & 0.0 [0.0, 0.0] & 0.0 [0.0, 0.0] & 46.0 [40.0, 55.7] & 35.1 [30.0, 41.2] \\
tickets & 163 & 0.0 [0.0, 0.0] & 0.0 [0.0, 0.0] & 56.4 [49.0, 64.4] & 52.8 [47.5, 56.9] \\
\bottomrule
\end{tabular}
\end{table}


\subsection{Step-level detection with intervals}\label{app:detection}

\begin{table}[!ht]
\centering\scriptsize
\setlength{\tabcolsep}{4pt}
\caption{Step-level detection with 95\% cluster-bootstrap intervals over task templates (\DetectionClusters{} clusters for the instrumented arms). Precision, recall, F1 and type accuracy are percentages; AUC is a proportion. The intervals are wide because the cluster bootstrap resamples whole templates, and three templates cannot pin a step-level rate to better than a few points. They describe the \emph{published} run: \texttt{scripts/detection\_intervals.py} replays the detector arm and aborts unless every point estimate reproduces \texttt{report\_main.json} exactly.}
\label{tab:detection_ci}
\begin{adjustbox}{max width=\textwidth}
\begin{tabular}{@{}llccccccc@{}}
\toprule
Detector & Monitor & Precision & Recall & F1 & AUC & Type & Steps & Cl. \\
\midrule
Rules & instr. & 66.9 [66.2, 68.1] & 41.9 [40.1, 44.8] & 51.5 [50.0, 53.4] & 0.772 [0.749, 0.789] & 34.6 [33.6, 36.2] & 24,175 & 3 \\
Rules & blackbox & 86.7 [85.7, 87.7] & 32.0 [31.7, 32.3] & 46.7 [46.5, 46.9] & 0.787 [0.759, 0.822] & 22.1 [21.7, 22.6] & 1,094 & 2 \\
Supervised & instr. & 48.8 [42.4, 54.5] & 86.2 [84.7, 89.1] & 62.3 [56.5, 67.6] & 0.762 [0.753, 0.770] & 51.6 [49.3, 53.8] & 11,193 & 3 \\
Supervised & blackbox & 46.0 [43.7, 48.7] & 81.0 [80.5, 81.5] & 58.7 [56.7, 60.9] & 0.750 [0.747, 0.752] & 39.9 [37.4, 42.7] & 1,060 & 2 \\
\bottomrule
\end{tabular}
\end{adjustbox}
\end{table}


\begin{table}[!ht]
\centering\scriptsize
\setlength{\tabcolsep}{4pt}
\caption{Step-level detection by split, per detector and monitor regime. Detection neither improves nor collapses out of distribution, which is what a taxonomy-defined monitor should show: the rules encode the signal, not the fault schedule.}
\label{tab:appendix_split_det}
\begin{adjustbox}{max width=\textwidth}
\begin{tabular}{@{}llrcc@{}}
\toprule
Arm & Split & Steps & AUC & Firing rate (\%) \\
\midrule
rules / instrumented & train & 10,186 & 0.755 & 20.4 \\
rules / instrumented & dev & 2,672 & 0.758 & 27.1 \\
rules / instrumented & test-id & 2,703 & 0.779 & 27.3 \\
rules / instrumented & test-ood & 8,614 & 0.789 & 21.5 \\
rules / blackbox & test-id & 491 & 0.822 & 13.2 \\
rules / blackbox & test-ood & 603 & 0.759 & 11.6 \\
supervised / instrumented & train & 4,170 & 0.745 & 67.8 \\
supervised / instrumented & dev & 1,426 & 0.774 & 58.7 \\
supervised / instrumented & test-id & 1,324 & 0.762 & 65.1 \\
supervised / instrumented & test-ood & 4,273 & 0.770 & 63.7 \\
supervised / blackbox & test-id & 487 & 0.752 & 61.2 \\
supervised / blackbox & test-ood & 573 & 0.747 & 61.1 \\
\bottomrule
\end{tabular}
\end{adjustbox}
\end{table}


\section{Paired significance tests}\label{app:tests}

The main text reports cluster-bootstrap intervals, which answer ``how large is the
difference''. This appendix adds the complementary question---``is a difference present at
all''---with an exact McNemar test on the aligned episode-level outcomes of two policies,
followed by a Holm--Bonferroni step-down correction within each family of tests. Only
discordant pairs carry information: $b$ is the number of episodes the first system recovered
and the second did not, $c$ the reverse. The two-sided $p$-value is an exact binomial test on
$b+c$ pairs at $p=0.5$; no chi-square approximation is used, and no $t$- or Wilcoxon test is
appropriate because the outcome is a paired bit.

Three families are corrected separately: the \HolmPolicyTests{} policy pairs, the
\HolmClassTests{} per-class rule-versus-ranker contrasts and the \HolmFamilyTests{}
per-family contrasts. Correcting across all of them at once would test a question nobody
asked. The LLM-controller arm is a case of its own: its 240 pairs are a different denominator,
drawn from the same grid but scored against a different controller, so it gets a single
pre-specified test rather than membership of a family. The canonical mapping recovers
$b=\LLMMcNemarB$ of those pairs and the LLM $c=\LLMMcNemarC$, an exact McNemar $p=\LLMMcNemarP$;
no correction applies because there is one test and it was not chosen after seeing the outcome.
\texttt{scripts/llm\_paired\_test.py} recomputes it from the released export, aborting unless the
joined outcomes reproduce \texttt{results/llm\_agent.json} exactly.

\begin{table}[!ht]
\centering\scriptsize
\setlength{\tabcolsep}{4pt}
\caption{Exact McNemar tests over the 861 harmful pairs: all six policy contrasts, Holm-corrected across the six. $b$ counts episodes the first system recovered and the second did not, $c$ the reverse. The canonical mapping beats the learned ranker: $b=104$ against $c=22$.}
\label{tab:appendix_tests_policy}
\begin{tabular}{@{}lrrrcc@{}}
\toprule
Comparison & Pairs & $b$ & $c$ & $p$ & $p_{\text{Holm}}$ \\
\midrule
No recovery vs Retry only & 861 & 0 & 0 & 1.000 & 1.000 \\
No recovery vs Rule-based & 861 & 0 & 524 & \ensuremath{<}0.001 & \ensuremath{<}0.001 \\
No recovery vs Model-based & 861 & 0 & 442 & \ensuremath{<}0.001 & \ensuremath{<}0.001 \\
Retry only vs Rule-based & 861 & 0 & 524 & \ensuremath{<}0.001 & \ensuremath{<}0.001 \\
Retry only vs Model-based & 861 & 0 & 442 & \ensuremath{<}0.001 & \ensuremath{<}0.001 \\
Rule-based vs Model-based & 861 & 104 & 22 & \ensuremath{<}0.001 & \ensuremath{<}0.001 \\
\bottomrule
\end{tabular}
\end{table}


\begin{table}[!ht]
\centering\scriptsize
\setlength{\tabcolsep}{4pt}
\caption{Per-class McNemar tests, rule-based versus model-based, Holm-corrected across the seven classes (significant in 2 of 7). Where the two policies differ at all, they differ in favour of the canonical mapping; the remaining classes agree closely, which is the claim the class table makes: repair depends on the operator, not the selector.}
\label{tab:appendix_tests_class}
\begin{adjustbox}{max width=\textwidth}
\begin{tabular}{@{}lrrrcc@{}}
\toprule
Failure class & Pairs & $b$ (rules) & $c$ (ranker) & $p$ & $p_{\text{Holm}}$ \\
\midrule
invalid\_arguments & 169 & 38 & 0 & \ensuremath{<}0.001 & \ensuremath{<}0.001 \\
stale\_state & 76 & 17 & 11 & 0.345 & 1.000 \\
contradictory\_evidence & 158 & 0 & 0 & 1.000 & 1.000 \\
planning\_error & 110 & 0 & 1 & 1.000 & 1.000 \\
communication\_failure & 45 & 0 & 0 & 1.000 & 1.000 \\
objective\_drift & 134 & 49 & 10 & \ensuremath{<}0.001 & \ensuremath{<}0.001 \\
unrecoverable\_external\_error & 169 & 0 & 0 & 1.000 & 1.000 \\
\bottomrule
\end{tabular}
\end{adjustbox}
\end{table}


\begin{table}[!ht]
\centering\scriptsize
\setlength{\tabcolsep}{4pt}
\caption{Per-family McNemar tests, rule-based versus model-based, Holm-corrected across the four families (significant in 3 of 4).}
\label{tab:appendix_tests_family}
\begin{tabular}{@{}lrrrcc@{}}
\toprule
Family & Pairs & $b$ (rules) & $c$ (ranker) & $p$ & $p_{\text{Holm}}$ \\
\midrule
inventory & 217 & 12 & 1 & 0.003 & 0.007 \\
migration & 196 & 34 & 0 & \ensuremath{<}0.001 & \ensuremath{<}0.001 \\
procurement & 285 & 42 & 11 & \ensuremath{<}0.001 & \ensuremath{<}0.001 \\
tickets & 163 & 16 & 10 & 0.327 & 0.327 \\
\bottomrule
\end{tabular}
\end{table}


\section{Grid, hyper-parameters and cost model}\label{app:grid}

\textbf{Grid.} \NumPairs{} (task, fault) pairs from four families $\times$ three templates
$\times$ five configurations $\times$ three seeds, of which \NumHarmful{} are harmful after
the screen; \NumEpisodes{} episodes are evaluated in total, including the natural-fault,
long-horizon, blackbox and ablation slices. Templates A/B and configurations
\texttt{cfg1}--\texttt{cfg3} are \emph{train}; \texttt{cfg5} is \emph{dev}; \texttt{cfg6} is an
unseen configuration (\emph{test-id}); template C is entirely held out (\emph{test-ood}).
Fault schedules are drawn per task from a seeded RNG, so memorising a schedule is not a
viable strategy.

\textbf{Detectors.} The rule monitor is 13 hand-written rules over the 19 monitor features; it
has no fitted parameters. The supervised detector is a multinomial logistic regression (8
classes including \emph{none}, L2 $=2\times10^{-3}$, learning rate $0.15$, 900 epochs, class
balancing) fit on \DetTrainRows{} training steps; its decision threshold ($\tau=\SupThreshold$)
is swept on \emph{dev} to maximise F1 and never on a reporting split. The recovery ranker is a
set of per-operator logistic value models fit on deliberate decisions from the rule-based
policy on \emph{train}, with operators seen fewer than 25 times excluded from the candidate set.

\textbf{Intervals.} 95\% cluster bootstrap resampling whole task templates (2000 resamples for
the reported tables), a paired bootstrap for policy differences, and Wilson intervals for
single proportions. Because the grid has only three templates, template-clustered intervals are
coarse by construction; they are reported rather than replaced by episode-level intervals,
which would pretend episodes are independent draws.

\textbf{Cost.} USD comes from the fixed price sheet \texttt{arb-price-2026.11} with
deterministic per-step jitter seeded by the episode id, so a cost figure is reproducible
rather than sampled. The sheet is a transparent normalisation of token, tool-call and latency
prices, not a market forecast.

\begin{sloppypar}
\textbf{Reproduction.} \texttt{python scripts/run\_experiments.py --preset main} writes the
report and the episode table; \texttt{python scripts/detection\_intervals.py} replays the
detector arm and writes the interval artifact, aborting if it cannot reproduce the report;
\texttt{python scripts/llm\_paired\_test.py} recomputes the LLM arm's paired test from the
released export, aborting unless it reproduces the scored aggregates;
\texttt{python paper/analysis.py} writes every table, figure and macro here; two
\texttt{pdflatex} passes and one \texttt{bibtex} pass rebuild this PDF. The harness is tag
\texttt{arb-1.0.1}, and the dataset revision must be pinned with it for any comparison.
\end{sloppypar}

\section{The controller prompt}\label{app:prompt}

The LLM-controller arm replaces only the operator-selection step. The prompt below is emitted
verbatim from \texttt{arb/agents/llm.py}, so the no-oracle claim can be checked without reading
code: it carries the detector's posterior over classes, the monitor feature vector and the
operator menu, and it contains no field identifying the injected fault, its injection step,
the expected state or the grader's verdict. The step number it quotes is the step at which the
alarm fired, not the step at which the fault was injected.

% Auto-generated by paper/analysis.py -- do not edit.
{\footnotesize
\begin{verbatim}
You are the recovery controller of a tool-using agent. A monitor has raised an
    alarm after a tool call. Choose exactly one recovery operator from the
    menu: the one whose precondition matches the evidence. Prefer the least
    disruptive operator that addresses the cause; do not invent operators.
    Reply with JSON only: {"action": "<operator>", "confidence": <0-1>,
    "reason": "<10 words max>"}.

% ---- user message ----
Alarm at step 7. Retries used 0/2.
Failure posterior: stale_state=0.710, planning_error=0.120,
    objective_drift=0.090, contradictory_evidence=0.080
Monitor evidence: version_lag=2.00, stale_flag=1.00, _score=0.66,
    step_pressure=0.40
Operators:
- continue: alarm is a false positive; keep executing the plan unchanged
- retry: re-issue the identical call once (transient transport error only)
- revalidate_args: repair the call's arguments against the tool schema, then
    re-issue
- refresh_state: re-read the object being mutated so the write uses current
    state
- re_read_authoritative: re-read the authoritative source and adopt its value
    over other observations
- replan: discard the remaining steps and recompute the plan from current
    state
- resync_message: re-send or re-read the inter-agent message so the recipient
    acts on the sender's plan
- reanchor_objective: restore the stated objective and recompute the plan
    under it
- escalate_abort: halt cleanly with no further mutations and escalate to a
    human
Choose one operator.
\end{verbatim}
}

